% !PRINTED 478
\weberpage{478}

Um die charakteristischen Eigenschaften der Galois'schen Gruppe zu finden, haben wir nur die Sätze \S.\,146, 1. bis 4. etwas anders auszudrücken. Wir sagen von einer Function der $m$ Wurzeln, die bei irgend einer Permutation ungeändert bleibt, sie \gesperrt{gestatte diese Permutation}, und erhalten so:

a) \gesperrt{Jede rationale Gleichung in $\Omega$, die zwischen den $m$ Wurzeln von $F(x)$ besteht, bleibt richtig, wenn die Wurzeln irgend einer Permutation der Galois'schen Gruppe unterworfen werden.}

b) \gesperrt{Jede rationale Function in $\Omega$ von den $m$ Wurzeln von $F(x)$, die sämmtliche Permutationen der Galois'schen Gruppe gestattet, ist eine Zahl in $\Omega$.}

Denn drücken wir die Wurzeln von $F(x)$ als rationale Functionen von $\varrho$ aus, so geht eine Function der Wurzeln $\alpha$ in eine Function $\varphi(\varrho)$ über. Hat man zuvor eine Permutation der Galois'schen Gruppe ausgeführt, so erhält man eine der conjugirten Zahlen $\varphi(\varrho_a)$, die aus $\varphi(\varrho)$ durch eine Substitution $\sigma_a = (\varrho, \varrho_a)$ entsteht. Ist nun $\varphi(\varrho) = 0$, so sind nach dem Satz \S.\,146, 1. auch alle conjugirten $\varphi(\varrho_a) = 0$, womit a) bewiesen ist; und sind die conjugirten Grössen $\varphi(\varrho_a)$ alle einander gleich, so ist ihr gemeinsamer Werth nach \S.\,146, 4. in $\Omega$ enthalten, wodurch b) bewiesen ist.

Zu a) b) kommt noch als Drittes:

c) \gesperrt{Wenn irgend eine Permutation $\pi$ der Wurzeln von $F(x)$ auf alle rationalen Gleichungen in $\Omega$, die zwischen den Wurzeln bestehen, anwendbar ist, so gehört $\pi$ der Galois'schen Gruppe an, und die Galois'sche Gruppe kann daher auch erklärt werden als der Inbegriff aller Permutationen, die auf sämmtliche rationale Gleichungen zwischen den Wurzeln anwendbar sind.}

Denn nach \S.\,143 kann man $\varrho$ als rationale (z.\,B. lineare) Function der $m$ Grössen $\alpha, \alpha_1, \alpha_2 \ldots \alpha_{m-1}$ so annehmen, dass alle $\Pi(m)$ Werthe, die man durch die $\Pi(m)$ Permutationen der $\alpha$ daraus erhält, von einander verschieden sind. Ist dann $g(t) = 0$ die Galois'sche Resolvente von $F(x) = 0$, deren Wurzel dieses $\varrho$ ist, so ist
\[
g(\varrho) = 0.
\]
